\documentclass{article}


\usepackage{lmodern}
\usepackage{xcolor}
\usepackage[utf8]{inputenc}
\usepackage[T1]{fontenc}
\usepackage[francais]{babel}
\usepackage[top=1.7cm, bottom=1.7cm, left=1.7cm, right=1.7cm]{geometry}
\usepackage{verbatim}
\usepackage{tikz} %Vectoriel
\usepackage{listings}
\usepackage{fancyhdr}
\usepackage{multido}
\usepackage{amssymb}
\usepackage{float}

\newcommand{\titre}{Les emprunts et les prêts}
\newcommand{\numTD}{1}

\newcommand{\module}{Création d'entreprise}
\newcommand{\sigle}{ce}

\newcommand{\semestre}{4}

\definecolor{gris1}{gray}{0.40}
\definecolor{gris2}{gray}{0.55}
\definecolor{gris3}{gray}{0.65}
\definecolor{gris4}{gray}{0.50}
\definecolor{vert}{rgb}{0,0.4,0}
\definecolor{violet}{rgb}{0.65, 0.2, 0.65}
\definecolor{bleu1}{rgb}{0,0,0.8}
\definecolor{bleu2}{rgb}{0,0.2,0.6}
\definecolor{bleu3}{rgb}{0,0.2,0.2}
\definecolor{rouge}{HTML}{F93928}


\lstdefinelanguage{algo}{%
   morekeywords={%
    %%% couleur 1
		importer, programme, glossaire, fonction, procedure, constante, type, 
	%%% IMPORT & Co.
		si, sinon, alors, fin, tantque, debut, faire, lorsque, fin lorsque, 
		declenche, declencher, enregistrement, tableau, retourne, retourner, =, pour, a,
		/=, <, >, traite,exception, 
	%%% types 
		Entier, Reel, Booleen, Caractere, Réél, Booléen, Caractère,
	%%% types 
		entree, maj, sortie,entrée,
	%%% types 
		et, ou, non,
	},
  sensitive=true,
  morecomment=[l]{--},
  morestring=[b]',
}

\lstset{language=algo,
    %%% BOUCLE, TEST & Co.
      emph={importer, programme, glossaire, fonction, procedure, constante, type},
      emphstyle=\color{bleu2},
    %%% IMPORT & Co.  
	emph={[2]
		si, sinon, alors, fin , tantque, debut, faire, lorsque, fin lorsque, 
		declencher, retourner, et, ou, non,enregistrement, retourner, retourne, 
		tableau, /=, <, =, >, traite,exception, pour, a
	},
      emphstyle=[2]\color{bleu1},
    %%% FONCTIONS NUMERIQUES
      emph={[3]Entier, Reel, Booleen, Caractere, Booléen, Réél, Caractère},
      emphstyle=[3]\color{gris1},
    %%% FONCTIONS NUMERIQUES
      emph={[4]entree, maj, sortie, entrée},	
      emphstyle=[4]\color{gris1},
}
\lstdefinelanguage{wl}{%
   morekeywords={%
    %%% couleur 1
		importer, programme, glossaire, fonction, procedure, constante, type, 
	%%% IMPORT & Co.
		si, sinon, alors, fin, TANTQUE, tantque, FIN, PROCEDURE, debut, faire, lorsque, 
		fin lorsque, declenche, declencher, enregistrement, tableau, retourne, retourner, =, 
		/=, <, >, traite,exception, 
	%%% types 
		Entier, Reel, Booleen, Caractere, Réél, Booléen, Caractère,
	%%% types 
		entree, maj, sortie,entrée,
	%%% types 
		et, ou, non,
	},
  sensitive=true,
  morecomment=[l]{//},
  morestring=[b]',
}

\lstset{language=wl,
    %%% BOUCLE, TEST & Co.
      emph={importer, programme, glossaire, fonction, procedure, constante, type},
      emphstyle=\color{bleu2},
    %%% IMPORT & Co.  
	emph={[2]
		si, sinon, alors, fin , tantque, debut, faire, lorsque, fin lorsque, 
		declencher, retourner, et, ou, non,enregistrement, retourner, retourne, 
		tableau, /=, <, =, >, traite,exception
	},
      emphstyle=[2]\color{bleu1},
    %%% FONCTIONS NUMERIQUES
      emph={[3]Entier, Reel, Booleen, Caractere, Booléen, Réél, Caractère},
      emphstyle=[3]\color{gris1},
    %%% FONCTIONS NUMERIQUES
      emph={[4]entree, maj, sortie, entrée},	
      emphstyle=[4]\color{gris1},
}
\lstdefinelanguage{css}{%
   morekeywords={%
    %%% couleur 1
		background, image, repeat, position, index, color, border, font, 
		size, url, family, style, variant, weight, letter, spacing, line, 
		height, text, decoration, align, indent, transform, shadow, 
		background, image, repeat, position, index, color, border, font, 
		size, url, family, style, variant, weight, letter, spacing, line, 
		height, text, decoration, align, indent, transform, shadow, 
		vertical, align, white, space, word, spacing,attachment, width, 
		max, min, margin, padding, clip, direction, display, overflow,
		visibility, clear, float, top, right, bottom, left, list, type, 
		collapse, side, empty, cells, table, layout, cursor, marks, page, break,
		before, after, inside, orphans, windows, azimuth, after, before, cue, 
		elevation, pause, play, during, pitch, range, richness, spek, header, 
		numeral, punctuation, rate, stress, voice, volume,
	%%% types 
		left, right, bottom, top, none, center, solid, black, blue, red, green,
	},
  sensitive=true,
  sensitive=true,
  morecomment=[s]{/*}{*/},
  morestring=[b]',
}
\lstset{language=css,
    %%% BOUCLE, TEST & Co.
      emph={
		background, image, repeat, position, index, color, border, font, 
		size, url, family, style, variant, weight, letter, spacing, line, 
		height, text, decoration, align, indent, transform, shadow, 
		background, image, repeat, position, index, color, border, font, 
		size, url, family, style, variant, weight, letter, spacing, line, 
		height, text, decoration, align, indent, transform, shadow, 
		vertical, align, white, space, word, spacing,attachment, width, 
		max, min, margin, padding, clip, direction, display, overflow,
		visibility, clear, float, top, right, bottom, left, list, type, 
		collapse, side, empty, cells, table, layout, cursor, marks, page, break,
		before, after, inside, orphans, windows, azimuth, after, before, cue, 
		elevation, pause, play, during, pitch, range, richness, spek, header, 
		numeral, punctuation, rate, stress, voice, volume,
	  },
      emphstyle=\color{bleu2},
    %%% FONCTIONS NUMERIQUES
      emph={[3]
		left, right, bottom, top,none, solid, black, blue, green,
		  },
      emphstyle=[3]\color{bleu3},
    %%% FONCTIONS NUMERIQUES
}

\lstset{language=SQL,
    %%% BOUCLE, TEST & Co.
      emph={INSERT, UPDATE, DELETE, WHERE, SET, GROUP, BY, ORDER, REFERENCES},
      emphstyle=\color{bleu2},
    %%% IMPORT & Co.  
	emph={[2]
		if, end, begin, then, for, each, else, after, of, on, to
	},
      emphstyle=[2]\color{bleu1},
    %%% FONCTIONS NUMERIQUES
      emph={[3]Entier, Reel, Booleen, Caractere, Booléen, Réél, Caractère},
      emphstyle=[3]\color{gris1},
    %%% FONCTIONS NUMERIQUES
      emph={[4]entree, maj, sortie, entrée},	
      emphstyle=[4]\color{gris1},
}
\lstdefinelanguage{ARM}{%
   morekeywords={%
   ADD, SUB, MOV, MUL, RSB,CMP, BLS, BLE, B,BHI,LDR,
   BGE, RSBLT, BGT, BEQ, BNE,BLT,BHS,STR,STRB
	},
  sensitive=true,
  morecomment=[l]{@},
  morestring=[b]',
}

\lstset{ % general style for listings 
   numbers=left 
   , literate={é}{{\'e}}1 {è}{{\`e}}1 {à}{{\`a}}1 {ê}{{\^e}}1 {É}{{\'E}}1 {ô}{{\^o}}1 {€}{{\euro}}1{°}{{$^{\circ}$}}1 {ç}{ {c}}1 {ù}{u}1
	, extendedchars=\true
   , tabsize=2 
   , frame=l
   , framerule=1.1pt
   , linewidth=520px
   , breaklines=true 
   , basicstyle=\footnotesize\ttfamily 
   , numberstyle=\tiny\ttfamily 
   , framexleftmargin=0mm 
   , xleftmargin=0mm 
   , captionpos=b 
	, keywordstyle=\color{bleu2}
	, commentstyle=\color{vert}
	, stringstyle=\color{rouge}
	, showstringspaces=false
	, extendedchars=true
	, mathescape=true
} 
%	\lstlistoflistings
%	\addcontentsline{toc}{part}{List of code examples}
 %prise en charge du langage algo

\usepackage{ifthen}
\date{\today}

\chead{}
\rhead{TD\no\numTD}
\lhead{\titre}
%\makeindex

\lfoot{Université Paul Sabatier Toulouse III}
\rfoot{\sigle\semestre}
%\rfoot{}
\cfoot{--~~\thepage~~--}

\makeglossary
\makeatletter
\def\clap#1{\hbox to 0pt{\hss #1\hss}}%

\def\haut#1#2#3{%
	\hbox to \hsize{%
		\rlap{\vtop{\raggedright #1}
	}%
	\hss
	\clap{\vtop{\centering #2}
}%
\hss
\llap{\vtop{\raggedleft #3}}}}%
\def\bas#1#2#3{%
	\hbox to \hsize{%
		\rlap{\vbox{
			\raggedright #1
		}
	}%
	\hss \clap{\vbox{\centering #2}}%
	\hss
	\llap{\vbox{\raggedleft #3}}}
}%
\def\maketitle{%
	\thispagestyle{empty}{%
		\haut{}{\@blurb}{}
		%	
		%\vfill

		\begin{center}
			\vspace{-1.5cm}
			\usefont{OT1}{ptm}{m}{n}
			\huge \@numeroTD \@title
		\end{center}
		\par
		\hrule height 1pt
		\par
		\vspace{1cm}
		\bas{}{}{}
}%
}
\def\date#1{\def\@date{#1}}
\def\author#1{\def\@author{#1}}
\def\numeroTD#1{\def\@numeroTD{#1}}
\def\title#1{\def\@title{#1}}
\def\location#1{\def\@location{#1}}
\def\blurb#1{\def\@blurb{#1}}
\date{\today}
\newboolean{monBool}
\setboolean{monBool}{true}
\author{}
\title{}
\ifthenelse{\equal{\numTD}{}}{
\numeroTD{}
}
{
	\numeroTD{TD \no\numTD~--- }
}
\location{Amiens}\blurb{}
%\makeatother
\title{\titre}
\author{%Semestre \semestre
}

\location{Toulouse}
\blurb{%
\vspace{-35px}
\begin{flushleft}
	Université Paul Sabatier -- Toulouse III\\
	IUT A - Toulouse Rangueil\\
\end{flushleft}
\begin{flushright}
	\vspace{-45px}
	\Large \textbf \module \\
	\normalsize \textit \today\\
	Semestre \semestre
	\vspace{30px}
\end{flushright}
}%



%\title{Cours \\ \titre}
%\date{\today\\ Semestre \semestre}

%\lhead{Cours: \titre}
%\chead{}
%\rhead{\thepage}

%\lfoot{Université Paul Sabatier Toulouse III}
%\cfoot{\thepage}
%\rfoot{\sigle\semestre}

\pagestyle{fancy}

\newcommand{\remarque}[1]{
	\begin{center}
	\medskip
	\colorbox{remarque}{
		\begin{minipage}{0.85\textwidth}\medskip\includegraphics[height=10px]{images/remarque.png} #1 \medskip\end{minipage}
	}
	\medskip
	\end{center}
}

\newcounter{exemples}

\newenvironment{exemple}[1]{
   \vspace{-2mm}

\refstepcounter{exemples}
   \begin{center}
	\medskip
      \begin{minipage}{0.9\linewidth}
}{%
~
      \end{minipage}
   \end{center}~
   \vspace{-2mm}
}%

\newcommand{\captionExemple}[1]{
	\begin{center}{\bsc{Exemple} \thechapter.\arabic{exemples}~--~}#1\end{center}
}


\begin{document}
	\maketitle
	\section{Intérêts simples}
	\begin{verbatim}
       47500
---///---]------------------------------------------------]--///-------------------------->	
	   14/10/N                                         28/03/N+1
	\end{verbatim}
	\paragraph{Nombre de jours} $17+30+31+31+28+28 = 165$	
	\begin{eqnarray*}
		\frac{165}{360} \times 47500 \times 0.125 = 2721.35
	\end{eqnarray*}
	\remarque{
		\textbf{Intérêts simples} $Intérets = C \times Taux \times duree$\newline
		\textbf{Intérêts composés}	$C_n = C_0 \times (1+i)^n $
	}
	\section{Valeurs acquises à intérêts simples}
	\begin{eqnarray*}
		I &=& \frac{6 + 30 + 31 + 30 + 31 + 16}{360} \times 38000 \times (0.016 \times 12) = \textbf{2918.4}\\
		C_n &=& 38000+2918.4 = \textbf{40918.4}
	\end{eqnarray*}
	\section{Valeur acquise à intérêts composés}
	\begin{eqnarray*}
		C_n &=& C_0 \times (1+i)^n\\
		C_8 &=& 10000 \times (1 + 0.125)^{8} = \textbf{25657.85}
	\end{eqnarray*}
	\section{Capitalisation semestrielle}
		\begin{eqnarray*}
			C_n &=& C_0 \times (1+i)^n\\
			C_7 &=& 10000 \times (1 + 0.06)^{7\times2} = \textbf{22609.04}
		\end{eqnarray*}
	\section{Valeur actuelle à intérêts composés}

	\begin{eqnarray*}
		535897 &=& C_1 \times (1 + 0.10)^{8}\\
		C_1 &=& \frac{535897}{(1 + 0.10)^8} = \textbf{249999.91}  
	\end{eqnarray*}
	\section{Valeur actuelle d'un capital}
	On ramène toutes les dépenses en $0$ c'est-à-dire au 1$^{er}$ Janvier 2010, on se pose donc la question \textit{``Combien fallait il placer en $0$ pour pouvoir ensuite faire ces dépenses ?''}. \newline
	On choisira donc la solution, qui demandera le placement le plus faible.
	\paragraph{Solution A}
	$400000 - (3\%~de~400000)~en~0 = 388000$
	\paragraph{Solution B}
	$215000 \times (1.09)^0 + 215000 \times (1.09)^{-1} = 412247.7$
	\paragraph{Solution C}
	$150000 \times (1.09)^{-1} + 150000 \times (1.09)^{-2} + 150000 \times (1.09)^{-3} = 379694.18$
	\paragraph{Conclusion}
	On constate que la solution qui demande le moins d'argent en $0$ est la solution C.  C'est donc la plus intéressante pour l'entreprise.
	\section{Emprunt remboursé par amortissements constants}
	\begin{table}[H]
		\centering
		\begin{tabular}{|p{2cm}|p{3.6cm}|c|c|c|p{3.2cm}|}
			\hline
			\textbf{Échéance} & \textbf{Capital restant dû en début de période} & \textbf{Intérêts} & \textbf{Amortissements}& \textbf{Semestrialité} & \textbf{Capital restant dû en fin de période }\\
			\hline
			31/08/N & 100000& 10000&20000&30000&80000\\
			\hline
			28/02/N+1 & 80000&8000&20000&28000 & 60000\\
			\hline
			31/08/N+1 & 60000&6000&20000&26000&40000\\
			\hline
			28/02/N+2 & 40000& 4000&20000&24000&20000\\
			\hline
			31/08/N+2 &20000&2000&20000&22000 & 0\\ 
			\hline
		\end{tabular}
		\caption{Tableau d'amortissement}
	\end{table}
	\section{Emprunt remboursé par annuités constantes}
	\remarque{
			\textbf{Formule}:  

			$ a = V_0 \times \frac{i}{1-(1 \times i)^{-n}} $
	}
\begin{eqnarray*}
	a = 150\;000 \times \frac{0.10}{1-1.10^{-5}} = 39\;569.62 &\approx& 39\;570
\end{eqnarray*}
	\begin{table}[H]
		\centering
		\begin{tabular}{|p{2cm}|p{2.8cm}|c|c|c|p{2.8cm}|}
			\hline
			\textbf{Échéance} & \textbf{Capital restant dû en début } & \textbf{Intérêts annuels} & \textbf{Amortissements}& \textbf{Annuités} \\
			\hline
			1/3/N+1 & $200\;000$&$28\;000$& $40\;641$ & $68\;641$\\ 
			\hline
			1/3/N+2 & $159\;359$&$22\;310.26$& $46\;330.74$ & $68\;641$\\ 	
			\hline
			1/3/N+3 & $113\;028.26$&$15\;823.95$& $52\;817641.04$ & $68\;641$\\ 
			\hline
			1/3/N+4 & $60\;211.21$&$8\;429.57$& $60\;211.21$ & $68\;640.78$\\ 
			\hline
		\end{tabular}
		\caption{Tableau d'amortissements}
	\end{table}
	\section{Étude d'un financement}
	\section{Financement de l'investissement}
	\begin{table}[H]
		\centering
		\begin{tabular}{|p{2cm}|p{2.8cm}|c|c|c|p{2.8cm}|}
			\hline
			\textbf{Échéance} & \textbf{Capital restant dû en début } & \textbf{Intérêts annuels} & \textbf{Amortissements}& \textbf{Annuités} \\
			\hline
			01/01/2011 & $80\;000$		& $84\;000$		&	$17\;111$		&	$25\;511$\\
			\hline
			01/01/2012 & $62\;889$		& $6\;603.35$	&	$18\;907.65$	&	$25\;511$\\
			\hline
			01/01/2013 & $43\;981.35$	& $4\;618.04$	&	$20\;892.96$	&	$25\;511$\\
			\hline
			01/01/2014 & $23\;088.39$	& $2\;424.28$	&	$23\;088.39$	&	$25\;512.67$\\
			\hline
		\end{tabular}
		\caption{Tableau d'amortissements}
	\end{table}

\end{document}






